% SHARD-ID: CA-14-KOO-SALEUR-PROTOTYPE
% SHARD-TITLE: Koo-Saleur as the Prototype
% SHARD-SUMMARY: Uses the local lattice-fermion source to show a rigorous precedent for constructing continuum symmetry generators from lattice Hamiltonian-density modes.
% SHARD-SUMMARY: Records the Koo-Saleur pattern: Fourier modes of h_j plus commutator corrections approximate Virasoro generators.
% SHARD-SUMMARY: Keeps the original Koo-Saleur paper as an acquisition target rather than an uncited authority.
% SHARD-KEYWORDS: Koo-Saleur, Virasoro, Hamiltonian density, Fourier modes, commutator correction, OAR

\section{Koo--Saleur as the Prototype}
\label{sec:koo-saleur-prototype}

\subsection{Status, Sources, and Dependencies}

\textbf{Status: sourced precedent shard.}  The original Koo--Saleur paper is not
yet registered as a canonical source here.  This shard cites the local
lattice-fermion/OAR treatment, which states and proves a free-fermion
Koo--Saleur convergence result.

Dependencies: CA-10--CA-13.

% Source: references/text/CFTFromLatticeFermions.txt:384 -- KS convergence relies on correct scaling representation and vacuum subtraction.
% Source: references/text/CFTFromLatticeFermions.txt:395 -- transverse-field Ising even subalgebra and KS approximants converge.
% Source: references/text/CFTFromLatticeFermions.txt:2128 -- continuum H is integral of Hamiltonian density.
% Source: references/text/CFTFromLatticeFermions.txt:2133 -- h is a local quantum field.
% Source: references/text/CFTFromLatticeFermions.txt:2134 -- h is given by chiral and anti-chiral stress tensors.
% Source: references/text/CFTFromLatticeFermions.txt:2163 -- Fourier modes H_k are linear combinations of Virasoro generators.
% Source: references/text/CFTFromLatticeFermions.txt:2182 -- Koo-Saleur proposal introduces lattice Fourier modes of Hamiltonian density.
% Source: references/text/CFTFromLatticeFermions.txt:2197 -- KS approximants are lattice analogues of L_k and anti-chiral L_k.
% Source: references/text/CFTFromLatticeFermions.txt:2233 -- smeared KS approximants converge to smeared Virasoro generators.
% Source: references/text/CFTFromLatticeFermions.txt:2247 -- unitary exponentials and automorphic dynamics from smeared approximants.

\subsection{Continuum Side}

In a two-dimensional CFT on the circle, the Hamiltonian is expressed as an
integral of a Hamiltonian density,
\[
  H=\int h(x)\,dx.
\]
The source identifies \(h(x)\) as a local quantum field built from the chiral and
anti-chiral energy-momentum tensors \(T,\bar T\).  Its Fourier modes \(H_k\)
are linear combinations of the Virasoro generators \(L_k,\bar L_k\).

Thus in this special setting the desired spacetime symmetries are not abstract:
they are generated by modes of the stress tensor.

\subsection{Lattice Side}

The Koo--Saleur move is to imitate the continuum construction using the lattice
Hamiltonian density.  From a lattice density
\[
  h^{(N)}:\Lambda_N\to \mathcal A_N
\]
one forms Fourier modes \(H_k^{(N)}\).  The cited Definition 4.1 then combines
these modes with commutators against the lattice Hamiltonian mode \(H_0^{(N)}\)
to produce lattice analogues of \(L_k\) and \(\bar L_k\).

The important structural point for our programme is:
\[
  \text{lattice symmetry generator}
  =
  \text{weighted Hamiltonian density}
  +
  \text{commutator correction}.
\]
This is the same motif as CA-12, but now with Fourier weights and a Virasoro
target rather than a first moment and a Poincare translation target.

\subsection{Scaling Representation Matters}

The source explicitly warns that convergence of these approximants depends on
the correct scaling-limit representation and on subtracting ground-state
expectation values, corresponding to normal ordering.  Therefore the compiler
must not output only formal sums.  It must also output:
\[
\begin{array}{ll}
  &\text{state sequence},\\
  &\text{renormalisation/scaling maps},\\
  &\text{vacuum subtraction policy},\\
  &\text{operator topology target}.
\end{array}
\]

\subsection{What This Foreshadows}

Koo--Saleur is a prototype for the lattice-symmetry compiler:
\[
  \{h_j\} \longmapsto
  \{H_k^{(N)},L_k^{(N)},\bar L_k^{(N)}\}
  \longrightarrow
  \{L_k,\bar L_k\}.
\]
The first arrow is algebraic and finite-scale.  The second is analytic and
requires a scaling theorem.

\subsection{Boundary}

This shard does not import the original 1994 paper as authority.  It records it
as a future acquisition target.  It also does not assert that arbitrary anyon
chains satisfy the same convergence theorem; the cited source treats
free-fermion systems and related examples.
